\documentclass[12pt,letterpaper,oneside]{article}
\input{ITESCA/maestria-en-gestion-administrativa/plan-de-negocios-mga/formato-itesca-plan-negocios-generadas.tex}
\begin{document}
\portadaITESCA{Administración de recursos humanos: puestos y funciones}{6551}
\textbf{Borrador hipotético pendiente de confirmación.} La presente propuesta define los puestos y las funciones necesarios para operar AM Taller Autocentro. Su contenido no acredita contrataciones, relaciones laborales ni prestación efectiva de servicios profesionales. Se adopta el supuesto de que, durante la fase inicial, una sola persona podría concentrar varias responsabilidades; no obstante, las funciones se presentan por separado para delimitar responsabilidades, establecer controles y facilitar una eventual ampliación del proyecto.

\section{Estructura organizativa propuesta}

La organización se divide en cuatro áreas funcionales: administración, ventas y atención, operación técnica, y servicios profesionales de apoyo. Esta separación responde a las necesidades previstas de un taller, llantera y refaccionaria con agenda de citas, control de inventario y revisión de proveedores. La definición de puestos no implica que todos deban contratarse desde el inicio, pues su incorporación dependerá del volumen real de trabajo, la capacidad financiera y las obligaciones aplicables.

La coordinación general supervisaría la recepción de solicitudes, programación de trabajos, autorización de compras y control administrativo. La operación comprendería los servicios de mecánica y llantas, mientras que la gestión de refacciones atendería cotizaciones, compatibilidades e inventario. Los servicios contables, jurídicos, tecnológicos y de seguridad podrían contratarse externamente cuando se confirme su necesidad y alcance.

\section{Puestos y funciones}

\subsection{Responsable de administración y coordinación general}

\textbf{Función principal:} dirigir la operación administrativa y coordinar los recursos humanos, materiales y financieros del proyecto.

\textbf{Actividades propuestas:}
\begin{itemize}
    \item Establecer objetivos de operación y revisar su cumplimiento.
    \item Organizar horarios, cargas de trabajo y prioridades de servicio.
    \item Autorizar compras, pagos y solicitudes de cotización.
    \item Supervisar los registros de ingresos, egresos y documentación administrativa.
    \item Coordinar la selección, inducción y evaluación del personal cuando existan contrataciones.
    \item Verificar que los trabajos cuenten con autorización antes de ejecutarse.
    \item Revisar la relación con proveedores potenciales sin asumir contratos todavía no documentados.
    \item Atender incidentes y decisiones que excedan la autoridad de otros puestos.
\end{itemize}

Durante el inicio, este puesto podría acumular las funciones de caja, compras y supervisión. En ese supuesto, cada movimiento deberá conservar evidencia documental y someterse a revisiones periódicas para reducir errores y mantener un control básico de los recursos.

\subsection{Responsable de recepción, ventas y atención}

\textbf{Función principal:} recibir solicitudes y mantener la comunicación relacionada con citas, cotizaciones y seguimiento de servicios.

\textbf{Actividades propuestas:}
\begin{itemize}
    \item Registrar los datos indispensables de contacto y del vehículo.
    \item Administrar la agenda de citas conforme a la capacidad operativa confirmada.
    \item Recibir solicitudes de servicios de mecánica, llantas o refacciones.
    \item Elaborar cotizaciones preliminares sin presentar como disponibles las referencias de catálogo pendientes de precio, existencia o compatibilidad.
    \item Canalizar las consultas técnicas al personal competente.
    \item Solicitar autorización ante cualquier modificación de alcance o costo.
    \item Registrar pagos únicamente mediante el procedimiento administrativo establecido.
    \item Resguardar los datos proporcionados conforme a las disposiciones aplicables.
\end{itemize}

La persona responsable deberá evitar diagnósticos no confirmados y promesas sobre resultados, tiempos, garantías o disponibilidad. Su comunicación deberá ser clara, respetuosa y congruente con la identidad existente de AM Taller Autocentro.

\subsection{Técnico responsable de mecánica}

\textbf{Función principal:} revisar los vehículos y ejecutar los trabajos mecánicos previamente diagnosticados y autorizados.

\textbf{Actividades propuestas:}
\begin{itemize}
    \item Realizar una inspección inicial y documentar las condiciones observables.
    \item Proponer procedimientos de diagnóstico antes de sustituir componentes.
    \item Informar los insumos, herramientas y tiempo estimado requeridos.
    \item Ejecutar solamente las actividades incluidas en la orden autorizada.
    \item Registrar trabajos efectuados, piezas utilizadas y observaciones posteriores.
    \item Mantener orden, limpieza y condiciones seguras en el área de trabajo.
    \item Reportar daños, riesgos o necesidades adicionales antes de intervenir.
    \item Participar en la capacitación técnica y de seguridad que se determine.
\end{itemize}

La asignación de trabajos deberá corresponder con la experiencia y las competencias comprobables de la persona técnica. Las intervenciones que excedan su capacidad deberán detenerse o canalizarse con un profesional competente.

\subsection{Técnico responsable de llantas}

\textbf{Función principal:} atender los servicios relacionados con revisión, instalación y mantenimiento de llantas dentro de la capacidad técnica disponible.

\textbf{Actividades propuestas:}
\begin{itemize}
    \item Revisar el estado visible de llantas, válvulas y componentes relacionados.
    \item Confirmar medidas y especificaciones antes de solicitar una cotización.
    \item Ejecutar los procedimientos autorizados con herramientas adecuadas.
    \item Aplicar medidas de seguridad durante la elevación, el desmontaje y el montaje.
    \item Registrar los materiales utilizados y mantener identificados los bienes recibidos.
    \item Informar sobre condiciones que requieran una revisión especializada.
\end{itemize}

Este puesto no presupone que el proyecto cuente actualmente con maquinaria, inventario o personal especializado. Tales recursos permanecen sujetos a confirmación antes de ofrecer cualquier servicio.

\subsection{Responsable de refacciones, compras e inventario}

\textbf{Función principal:} controlar las solicitudes, entradas, salidas y resguardo de refacciones e insumos.

\textbf{Actividades propuestas:}
\begin{itemize}
    \item Solicitar cotizaciones a proveedores potenciales y comparar sus condiciones.
    \item Verificar números de parte, medidas y compatibilidades con apoyo técnico.
    \item Registrar entradas y salidas asociándolas con una orden de servicio.
    \item Identificar diferencias entre existencias físicas y registros.
    \item Separar materiales pendientes de revisión, devolución o reclamación.
    \item Proponer reposiciones con base en información real, sin asumir una demanda validada.
    \item Conservar comprobantes de compra y documentos de recepción.
\end{itemize}

Si una sola persona realiza compras e inventario, la administración deberá revisar cotizaciones, autorizaciones y recepciones para reducir errores y posibles conflictos de interés.

\subsection{Servicios profesionales externos}

\textbf{Función principal:} proporcionar asesoría especializada sin integrarse necesariamente a la plantilla permanente.

Podrían considerarse servicios de contabilidad y cumplimiento fiscal, asesoría jurídica, seguridad e higiene, mantenimiento de equipos y soporte para la agenda y los registros. Cada contratación requeriría una cotización, un alcance definido, la identificación de responsables, condiciones de confidencialidad y entregables verificables. Su incorporación dependerá de la formalización del proyecto y de las obligaciones aplicables.

\section{Asignación y control de responsabilidades}

Cada solicitud deberá identificar responsables de recepción, autorización administrativa, ejecución técnica y control de insumos, aunque inicialmente varias funciones recaigan en una sola persona. El flujo propuesto es el siguiente:

\begin{enumerate}
    \item Recepción registra la solicitud y los datos indispensables.
    \item El área técnica revisa el vehículo y propone el alcance.
    \item Administración valida los recursos y prepara la cotización.
    \item La persona solicitante autoriza antes del inicio del trabajo.
    \item El personal técnico ejecuta y documenta el servicio.
    \item Inventario registra los insumos utilizados.
    \item Recepción comunica la conclusión y conserva la evidencia correspondiente.
\end{enumerate}

La evaluación interna podría considerar cumplimiento de procedimientos, calidad de los registros, cuidado de herramientas, atención respetuosa, puntualidad y reporte oportuno de incidentes. Posteriormente, estos criterios deberán convertirse en descripciones formales de puesto e indicadores verificables.

\section{Datos requeridos para la versión definitiva}

Para convertir este borrador en un trabajo real se necesita confirmar la cantidad de personas participantes, las funciones actualmente realizadas, horarios, competencias técnicas, capacidad instalada, equipo disponible, modalidad de contratación, presupuesto de nómina y necesidad de servicios externos. También permanecen pendientes el formato, la rúbrica y el peso individual de la actividad, así como la validación de las obligaciones laborales, fiscales, de protección de datos y de seguridad aplicables a AM Taller Autocentro.
\clearpage
\printbibliography[title={Fuentes de consulta}]
\end{document}
